% =====================================================================
% 7. Code distribution and reproduction
% =====================================================================
\section{Code distribution and reproduction}\label{sec:code}

\subsection{Repository}

All code, run logs, figures, tables and the sources of the four reports are in the public repository
\url{https://github.com/orianafe16/ptbxl-DL}:

\begin{lstlisting}[language={},numbers=none,caption={Structure of the repository.},label={lst:tree}]
src/          data, models (src/models/), training loop, metrics, one *_experiments.py per block
tests/        pytest suite, one file per block (synthetic data; the real cache when present)
notebooks/    EDA and zero-shot notebooks (they call the functions of src/)
results/logs/ one JSON per run and one summary per block (checkpoints and caches not committed)
report/       numbers.tex (generated), runs.json, published.json, tables/, figures/,
              common/ (shared LaTeX), one main file per report
docs/         DECISIONS.md, PLAN_DL.md, AI_USAGE_LOG.md, blocks/ (one note per block),
              literature/ (prompts, checked quotes, learning-strategies material)
\end{lstlisting}

The data are not committed: the README explains how to download PTB-XL 1.0.3 from PhysioNet, and
\texttt{python -m src.prepare\_data} checks the files against the dataset's checksums and rebuilds the
cache. Checkpoints and smoke-run outputs are ignored by git.

\subsection{Reproducing the results}

On a Windows laptop with an NVIDIA GPU (the README also covers the CPU case):
\begin{lstlisting}[language={},numbers=none,caption={Commands to reproduce the project, from the folder \texttt{ptbxl-dl}.},label={lst:repro}]
py -3.12 -m venv .venv ; .venv\Scripts\activate
pip install torch --index-url https://download.pytorch.org/whl/cu126
pip install -r requirements.txt
python -m src.check_env ; python -m pytest -q tests
python -m src.prepare_data                       # checksums + cache
python -m src.eda                                # EDA: eda.json and figures
python -m src.train --model trivial --final      # trivial reference
python -m src.dense_experiments all              # likewise cnn_, rnn_, transformer_
python -m src.generative_experiments all         # generators (no test fold)
python -m src.transfer_experiments run ; python -m src.transfer_experiments report
python -m src.leads_experiments all ; python -m src.zero_shot run
python -m src.comparison report                  # final table and figures
python report/make_numbers.py                    # report/numbers.tex
python report/common/make_floats.py              # shared figure/table floats
python report/common/make_contributions.py       # contribution tables from PLAN_DL.md
python report/common/make_code_excerpts.py       # code listings of this report
\end{lstlisting}

The order matters: each block compares its final models with those of earlier blocks (Dense, then CNN,
RNN and Transformer), the generators must exist before the transfer experiments, and the comparison comes
last. Every block also has a smoke mode (\texttt{all --fast}) and, where runs are long, an \texttt{estimate}
step. Reproduction has three levels. Without data or GPU, a clean clone runs the tests on synthetic data
and rebuilds \file{numbers.tex} from the committed logs (we checked that the macros come out identical),
which is all the reports need. With the data and the checkpoints (kept outside git and backed up
separately), the \texttt{report} steps rebuild every table and figure, including those that load a model.
Retraining repeats the runs, with small GPU differences because the fully deterministic mode of PyTorch is
not enabled. The reports compile from \texttt{report/} (on Overleaf, the folder \texttt{report/} uploaded as a
project with \texttt{oriana\_mllb.tex} as main file).

\subsection{Tests}

Each block has its tests, run before every long execution: data (shapes, no patient in two splits, class
counts against PhysioNet, normalisation statistics from training records only), metrics on toy cases,
the loop over-fitting a small batch, receptive field and class activation maps of the CNN, attention
against \texttt{nn.MultiheadAttention} with and without the causal mask, causality of the generators,
the tokenizer, the frozen backbone of the probe, and every pipeline end to end on synthetic data.
Regression tests were added for each code error found (Appendix~\ref{app:ai-code}).

\subsection{Final results and traceability}

% Shared float (common/make_floats.py). Table: tables/comparison_final.tex
\begin{table}[htbp]
\centering
\caption[Final comparison (test)]{Final comparison on the test fold. Mean $\pm$ SD over seeds; interval: 95\,\% patient bootstrap of the first seed. Published rows: macro AUC of the superclass task in Table~I of \cite{strodthoff2021deep}, obtained in a different setting (described in the text).}
\label{tab:comparison_final}
\small
\begin{adjustbox}{max width=\textwidth}
% Generated by src/comparison.py from results/logs/comparison_summary.json.
% Mean +- SD over seeds (one run: no SD); CI = 95% patient bootstrap of the first seed. Needs alias "comparison".
\begin{tabular}{lcrrccc}
\toprule
Model & Labels & Train.\ param. & Seeds & Test macro AUROC & First seed [95\% CI] & Test macro F1 \\
\midrule
\multicolumn{7}{l}{\textit{12 leads}} \\
\quad Trivial (lead mean and SD) & 100\% & \comparisonParamsTrivial{} & \comparisonNSeedsTrivial{} & \comparisonTestAurocMeanTrivial{} & \comparisonTestAurocFirstTrivial{} [\comparisonTestAurocLoFirstTrivial{}, \comparisonTestAurocHiFirstTrivial{}] & \comparisonTestFOneMeanTrivial{} \\
\quad Dense & 100\% & \comparisonParamsTwelveDense{} & \comparisonNSeedsTwelveDense{} & \comparisonTestAurocMeanTwelveDense{} $\pm$ \comparisonTestAurocSdTwelveDense{} & \comparisonTestAurocFirstTwelveDense{} [\comparisonTestAurocLoFirstTwelveDense{}, \comparisonTestAurocHiFirstTwelveDense{}] & \comparisonTestFOneMeanTwelveDense{} $\pm$ \comparisonTestFOneSdTwelveDense{} \\
\quad CNN & 100\% & \comparisonParamsTwelveCnn{} & \comparisonNSeedsTwelveCnn{} & \comparisonTestAurocMeanTwelveCnn{} $\pm$ \comparisonTestAurocSdTwelveCnn{} & \comparisonTestAurocFirstTwelveCnn{} [\comparisonTestAurocLoFirstTwelveCnn{}, \comparisonTestAurocHiFirstTwelveCnn{}] & \comparisonTestFOneMeanTwelveCnn{} $\pm$ \comparisonTestFOneSdTwelveCnn{} \\
\quad RNN & 100\% & \comparisonParamsTwelveRnn{} & \comparisonNSeedsTwelveRnn{} & \comparisonTestAurocMeanTwelveRnn{} $\pm$ \comparisonTestAurocSdTwelveRnn{} & \comparisonTestAurocFirstTwelveRnn{} [\comparisonTestAurocLoFirstTwelveRnn{}, \comparisonTestAurocHiFirstTwelveRnn{}] & \comparisonTestFOneMeanTwelveRnn{} $\pm$ \comparisonTestFOneSdTwelveRnn{} \\
\quad RNN + attention & 100\% & \comparisonParamsTwelveRnnattn{} & \comparisonNSeedsTwelveRnnattn{} & \comparisonTestAurocMeanTwelveRnnattn{} $\pm$ \comparisonTestAurocSdTwelveRnnattn{} & \comparisonTestAurocFirstTwelveRnnattn{} [\comparisonTestAurocLoFirstTwelveRnnattn{}, \comparisonTestAurocHiFirstTwelveRnnattn{}] & \comparisonTestFOneMeanTwelveRnnattn{} $\pm$ \comparisonTestFOneSdTwelveRnnattn{} \\
\quad Transformer & 100\% & \comparisonParamsTwelveTransformer{} & \comparisonNSeedsTwelveTransformer{} & \comparisonTestAurocMeanTwelveTransformer{} $\pm$ \comparisonTestAurocSdTwelveTransformer{} & \comparisonTestAurocFirstTwelveTransformer{} [\comparisonTestAurocLoFirstTwelveTransformer{}, \comparisonTestAurocHiFirstTwelveTransformer{}] & \comparisonTestFOneMeanTwelveTransformer{} $\pm$ \comparisonTestFOneSdTwelveTransformer{} \\
\addlinespace
\multicolumn{7}{l}{\textit{Lead II only (T14)}} \\
\quad Dense & 100\% & \comparisonParamsOneleadDense{} & \comparisonNSeedsOneleadDense{} & \comparisonTestAurocMeanOneleadDense{} $\pm$ \comparisonTestAurocSdOneleadDense{} & \comparisonTestAurocFirstOneleadDense{} [\comparisonTestAurocLoFirstOneleadDense{}, \comparisonTestAurocHiFirstOneleadDense{}] & \comparisonTestFOneMeanOneleadDense{} $\pm$ \comparisonTestFOneSdOneleadDense{} \\
\quad CNN & 100\% & \comparisonParamsOneleadCnn{} & \comparisonNSeedsOneleadCnn{} & \comparisonTestAurocMeanOneleadCnn{} $\pm$ \comparisonTestAurocSdOneleadCnn{} & \comparisonTestAurocFirstOneleadCnn{} [\comparisonTestAurocLoFirstOneleadCnn{}, \comparisonTestAurocHiFirstOneleadCnn{}] & \comparisonTestFOneMeanOneleadCnn{} $\pm$ \comparisonTestFOneSdOneleadCnn{} \\
\quad RNN & 100\% & \comparisonParamsOneleadRnn{} & \comparisonNSeedsOneleadRnn{} & \comparisonTestAurocMeanOneleadRnn{} $\pm$ \comparisonTestAurocSdOneleadRnn{} & \comparisonTestAurocFirstOneleadRnn{} [\comparisonTestAurocLoFirstOneleadRnn{}, \comparisonTestAurocHiFirstOneleadRnn{}] & \comparisonTestFOneMeanOneleadRnn{} $\pm$ \comparisonTestFOneSdOneleadRnn{} \\
\quad RNN + attention & 100\% & \comparisonParamsOneleadRnnattn{} & \comparisonNSeedsOneleadRnnattn{} & \comparisonTestAurocMeanOneleadRnnattn{} $\pm$ \comparisonTestAurocSdOneleadRnnattn{} & \comparisonTestAurocFirstOneleadRnnattn{} [\comparisonTestAurocLoFirstOneleadRnnattn{}, \comparisonTestAurocHiFirstOneleadRnnattn{}] & \comparisonTestFOneMeanOneleadRnnattn{} $\pm$ \comparisonTestFOneSdOneleadRnnattn{} \\
\quad Transformer & 100\% & \comparisonParamsOneleadTransformer{} & \comparisonNSeedsOneleadTransformer{} & \comparisonTestAurocMeanOneleadTransformer{} $\pm$ \comparisonTestAurocSdOneleadTransformer{} & \comparisonTestAurocFirstOneleadTransformer{} [\comparisonTestAurocLoFirstOneleadTransformer{}, \comparisonTestAurocHiFirstOneleadTransformer{}] & \comparisonTestFOneMeanOneleadTransformer{} $\pm$ \comparisonTestFOneSdOneleadTransformer{} \\
\addlinespace
\multicolumn{7}{l}{\textit{Lead II, ECG-GPT transfer (T13)}} \\
\quad ECG-GPT, from scratch & 100\% & \comparisonParamsScratchFOneZeroZero{} & \comparisonNSeedsScratchFOneZeroZero{} & \comparisonTestAurocMeanScratchFOneZeroZero{} $\pm$ \comparisonTestAurocSdScratchFOneZeroZero{} & \comparisonTestAurocFirstScratchFOneZeroZero{} [\comparisonTestAurocLoFirstScratchFOneZeroZero{}, \comparisonTestAurocHiFirstScratchFOneZeroZero{}] & \comparisonTestFOneMeanScratchFOneZeroZero{} $\pm$ \comparisonTestFOneSdScratchFOneZeroZero{} \\
\quad ECG-GPT, linear probe (frozen) & 100\% & \comparisonParamsProbeFOneZeroZero{} & \comparisonNSeedsProbeFOneZeroZero{} & \comparisonTestAurocMeanProbeFOneZeroZero{} $\pm$ \comparisonTestAurocSdProbeFOneZeroZero{} & \comparisonTestAurocFirstProbeFOneZeroZero{} [\comparisonTestAurocLoFirstProbeFOneZeroZero{}, \comparisonTestAurocHiFirstProbeFOneZeroZero{}] & \comparisonTestFOneMeanProbeFOneZeroZero{} $\pm$ \comparisonTestFOneSdProbeFOneZeroZero{} \\
\quad ECG-GPT, fine-tuning & 100\% & \comparisonParamsFinetuneFOneZeroZero{} & \comparisonNSeedsFinetuneFOneZeroZero{} & \comparisonTestAurocMeanFinetuneFOneZeroZero{} $\pm$ \comparisonTestAurocSdFinetuneFOneZeroZero{} & \comparisonTestAurocFirstFinetuneFOneZeroZero{} [\comparisonTestAurocLoFirstFinetuneFOneZeroZero{}, \comparisonTestAurocHiFirstFinetuneFOneZeroZero{}] & \comparisonTestFOneMeanFinetuneFOneZeroZero{} $\pm$ \comparisonTestFOneSdFinetuneFOneZeroZero{} \\
\quad ECG-GPT, from scratch & 10\% & \comparisonParamsScratchFOneZero{} & \comparisonNSeedsScratchFOneZero{} & \comparisonTestAurocMeanScratchFOneZero{} $\pm$ \comparisonTestAurocSdScratchFOneZero{} & \comparisonTestAurocFirstScratchFOneZero{} [\comparisonTestAurocLoFirstScratchFOneZero{}, \comparisonTestAurocHiFirstScratchFOneZero{}] & \comparisonTestFOneMeanScratchFOneZero{} $\pm$ \comparisonTestFOneSdScratchFOneZero{} \\
\quad ECG-GPT, linear probe (frozen) & 10\% & \comparisonParamsProbeFOneZero{} & \comparisonNSeedsProbeFOneZero{} & \comparisonTestAurocMeanProbeFOneZero{} $\pm$ \comparisonTestAurocSdProbeFOneZero{} & \comparisonTestAurocFirstProbeFOneZero{} [\comparisonTestAurocLoFirstProbeFOneZero{}, \comparisonTestAurocHiFirstProbeFOneZero{}] & \comparisonTestFOneMeanProbeFOneZero{} $\pm$ \comparisonTestFOneSdProbeFOneZero{} \\
\quad ECG-GPT, fine-tuning & 10\% & \comparisonParamsFinetuneFOneZero{} & \comparisonNSeedsFinetuneFOneZero{} & \comparisonTestAurocMeanFinetuneFOneZero{} $\pm$ \comparisonTestAurocSdFinetuneFOneZero{} & \comparisonTestAurocFirstFinetuneFOneZero{} [\comparisonTestAurocLoFirstFinetuneFOneZero{}, \comparisonTestAurocHiFirstFinetuneFOneZero{}] & \comparisonTestFOneMeanFinetuneFOneZero{} $\pm$ \comparisonTestFOneSdFinetuneFOneZero{} \\
\addlinespace
\multicolumn{7}{l}{\textit{Lead II, zero-shot (T24)}} \\
\quad CLIP ViT-B/32, zero-shot & 0\% & --- & \comparisonNSeedsZeroshot{} & \comparisonTestAurocMeanZeroshot{} & \comparisonTestAurocFirstZeroshot{} [\comparisonTestAurocLoFirstZeroshot{}, \comparisonTestAurocHiFirstZeroshot{}] & \comparisonTestFOneMeanZeroshot{} \\
\addlinespace
\multicolumn{7}{l}{\textit{Published (quoted, different settings: see text)}} \\
\quad resnet1d\_wang, 12 leads~\cite{strodthoff2021deep} & & & & \comparisonPublishedStrodthoffResnetOneDWang{} & & \\
\quad xresnet1d101, 12 leads~\cite{strodthoff2021deep} & & & & \comparisonPublishedStrodthoffXresnetOneDOneZeroOne{} & & \\
\quad lstm, 12 leads~\cite{strodthoff2021deep} & & & & \comparisonPublishedStrodthoffLstm{} & & \\
\quad fcn\_wang, 12 leads~\cite{strodthoff2021deep} & & & & \comparisonPublishedStrodthoffFcnWang{} & & \\
\quad inception1d, 12 leads~\cite{strodthoff2021deep} & & & & \comparisonPublishedStrodthoffInceptionOneD{} & & \\
\quad lstm\_bidir, 12 leads~\cite{strodthoff2021deep} & & & & \comparisonPublishedStrodthoffLstmBidir{} & & \\
\quad Wavelet+NN, 12 leads~\cite{strodthoff2021deep} & & & & \comparisonPublishedStrodthoffWaveletNn{} & & \\
\quad ensemble, 12 leads~\cite{strodthoff2021deep} & & & & \comparisonPublishedStrodthoffEnsemble{} & & \\
\bottomrule
\end{tabular}

\end{adjustbox}
\end{table}


Table~\ref{tab:comparison_final} collects every final model; each value is a macro whose source is the
summary log of its block, and each summary lists the run logs it read. The runs of the learning-strategy
block were made one commit before their code was committed (the logs say so); we checked on the laptop
that every file under \texttt{src/} was the delivered version and committed it unchanged.
